\begin{center}\begin{tikzpicture}[font=\small,>=stealth,line join=round]
\draw[gray!40,rounded corners=3pt] (0,0) rectangle (14.2,9.4);
\node at (7.1,8.8) {例3：原点不在定义域，也可以关于原点对称};
\begin{scope}[shift={(3.6,4.9)},x=1.05cm,y=0.6cm]\draw[->,gray] (-2.6,0)--(2.8000000000000003,0) node[right]{$x$};\draw[->,gray] (0,-5.5)--(0,5.7) node[left]{$y$};\begin{scope}\clip (-2.6,-5.5) rectangle (2.6,5.5);\draw[blue,thick,domain=-2.6:-0.185,samples=120] plot(\x,{\x+1/\x});\draw[blue,thick,domain=0.185:2.6,samples=120] plot(\x,{\x+1/\x});\end{scope}\draw[orange!85!black,dashed] (-1,-2)--(1,2);\fill[orange!85!black] (-1,-2) circle(2pt) node[above left]{$(-1,-2)$};\fill[orange!85!black] (1,2) circle(2pt) node[below right]{$(1,2)$};\draw[red!65!black] (0,0) circle(2pt);\node[below right,red!65!black] at (0,0) {原点不在图上};\end{scope}\node[align=left,anchor=west] at (7.9,5.3) {$f(x)=x+1/x$\\[12pt]$D=\mathbb R\setminus\{0\}$ 对称\\[12pt]$f(-x)=-f(x)$\\[12pt]奇函数可以没有 $x=0$};
\node[text width=12.8cm,align=center] at (7.1,.65) {红圈只提示原点缺失；曲线在 $x=0$ 两侧趋向无穷，视窗裁切不是端点。};
\end{tikzpicture}\end{center}
